\documentclass[11pt]{article}
% Shared preamble for the MSH_Tensor_Network_Methods function references (compile with tectonic / XeLaTeX)
\usepackage{fontspec}
\usepackage{unicode-math}
\setmainfont{texgyrepagella}[Extension=.otf,UprightFont=*-regular,BoldFont=*-bold,ItalicFont=*-italic,BoldItalicFont=*-bolditalic]
\setmathfont{texgyrepagella-math.otf}
\setmonofont{DejaVu Sans Mono}[Scale=0.82]
\newfontfamily\symfb{DejaVu Sans}[Scale=0.82]
\usepackage{newunicodechar}
\newunicodechar{ℓ}{{\symfb ℓ}}
\newunicodechar{∥}{{\symfb ∥}}
\newunicodechar{≪}{{\symfb ≪}}
\usepackage[a4paper,margin=2.2cm]{geometry}
\usepackage{amsmath}
\usepackage{booktabs,tabularx,longtable,array}
\usepackage{graphicx}
\usepackage{xcolor}
\usepackage{float}
\usepackage{caption}
\captionsetup{font=small,labelfont=bf}
\usepackage{fvextra}
\fvset{breaklines,breakanywhere,fontsize=\small,frame=leftline,framerule=1.5pt,rulecolor=\color{gray!45},xleftmargin=6pt}
\usepackage{enumitem}
\setlist{nosep,leftmargin=*}
\usepackage[hidelinks]{hyperref}
\usepackage{titlesec}
\titleformat{\section}{\Large\bfseries}{\thesection}{0.6em}{}
\titleformat{\subsection}{\large\bfseries}{\thesubsection}{0.6em}{}
\titleformat{\subsubsection}{\normalsize\bfseries}{\thesubsubsection}{0.6em}{}
\setlength{\parskip}{0.45em}\setlength{\parindent}{0pt}
% one box per function: \fn{name}{signature} then a description
\newcommand{\fnname}[1]{\texttt{\textbf{#1}}}
\newenvironment{function}[1]{\par\medskip\noindent\textbf{\large\texttt{#1}}\par\nopagebreak}{\par\medskip}
\newcommand{\code}[1]{\texttt{#1}}

\usepackage{mathtools}
\sloppy\emergencystretch=3em
\usepackage{needspace}
\renewenvironment{function}[1]{\par\Needspace{7\baselineskip}\medskip\noindent{\raggedright\textbf{\large\texttt{\hyphenchar\font=-1 #1}}\par}\nopagebreak}{\par\medskip}
\renewcommand{\code}[1]{\texttt{\hyphenchar\font=-1 #1}}
\hypersetup{pdftitle={MSH 3-T-junction array, logical gates: function reference}}
\newcommand{\ket}[1]{|#1\rangle}
\newcommand{\bra}[1]{\langle#1|}

\begin{document}

\begin{center}
{\LARGE\bfseries MSH 3-T-junction array: logical gates}\\[0.5em]
{\Large Function reference}\\[1.2em]
\begin{tabular}{lp{8.6cm}}
\code{logical\_gates\_tn.ipynb} & Julia: reduced method from the tensor-network (TDVP) evolution\\
\code{logical\_gates\_python\_ed.ipynb} & Python: full Pfaffian method and reduced method, dense\\
\code{tjunctions\_model.py} & Python: shared model of the array (with \code{tjunctions\_python\_ed.ipynb})\\
\end{tabular}\\[1.2em]
Repository \code{MSH\_Tensor\_Network\_Methods}, branch \code{tjunctions-logical-gates}\\[0.3em]
30 October 2026
\end{center}

\vspace{1em}
\noindent This document explains every function defined in the two logical-gate notebooks, and the functions of
\code{tjunctions\_model.py} they use: what each computes, the method, inputs and outputs, and the numerical settings that matter.
All numbers are quoted from the saved outputs of the executed notebooks. The model itself (geometry, Hamiltonian, protocols, zero-mode
construction) is described in detail in \code{docs/msh\_tjunctions\_function\_reference.pdf}; here it is only summarized.

\tableofcontents
\clearpage

%======================================================================
\section{Overview}
%======================================================================

\subsection{What is computed}

The array of \code{msh\_tjunctions\_tn.ipynb} carries four Majorana zero modes $\gamma_1,\dots,\gamma_4$, paired into the positive-energy
fermion modes $\psi_A = (\gamma_1 + i\gamma_2)/\sqrt2$ (junction 1) and $\psi_B = (\gamma_3 + i\gamma_4)/\sqrt2$ (junction 3). The logical qubit
is the even-parity pair
\[
\ket0_L = \ket{00},\qquad \ket1_L = \ket{11},
\]
where $\ket{n_An_B}$ has mode $\psi_m$ occupied when $n_m = 1$ (the convention of the masks in the CZ, QFT and period-finding notebooks).
The two gate protocols of the T-junction notebooks are analysed after one and after two exchanges:

\begin{center}
\begin{tabular}{lllcc}
\toprule
Protocol & Exchanges & $T_1$ & One exchange & Two\\
\midrule
\code{exchange\_12} (cell 1, \code{rot\_Tz = 200}) & $\gamma_1\leftrightarrow\gamma_2$ (same fermion) & 2040 & $\sqrt{Z}{}^\dagger$ & $Z$\\
\code{exchange\_23} (cell 2, \code{rot\_Tx = 400}) & $\gamma_2\leftrightarrow\gamma_3$ (different fermions) & 3880 & $\sqrt{X}{}^\dagger$ & $X$\\
\bottomrule
\end{tabular}
\end{center}

The result for each is the \textbf{logical unitary} on $\{\ket0_L,\ket1_L\}$, with signs and relative phases, by two methods:

\begin{itemize}
\item \textbf{Full method} (\code{signed\_amplitudes}): the many-body amplitudes $\bra j\hat U\ket k$ of the Gaussian (BCS) states from the
single-particle propagator $S_T$ applied to all $2N$ occupied quasiparticles, via the generalized Wick theorem (a Pfaffian). Gives in
addition the probability that the bulk stays unexcited. Needs the dense $4N\times4N$ propagator.
\item \textbf{Reduced method} (\code{braid\_matrix}, \code{logical\_unitary}): only the zero-mode block, the real $4\times4$ Majorana matrix
$M_{ij} = \bra{\gamma_i(0)}S_T\ket{\gamma_j(0)}$ and its spin representation. Needs only the two evolved mode vectors, so it works with
MPS (and with sparse propagators); it cannot see bulk excitations.
\end{itemize}

\subsection{Workflow and inputs}

\begin{center}\code{msh\_tjunctions\_tn.ipynb} (TDVP caches) $\to$ \code{logical\_gates\_tn.ipynb} $\to$ \code{logical\_gates\_tn.csv} $\to$ \code{logical\_gates\_python\_ed.ipynb}\end{center}

\begin{center}
\begin{tabularx}{\textwidth}{lX}
\toprule
File & Role\\
\midrule
\code{msh\_tjunctions\_tdvp\_\{Z,X\}.jls} & TDVP results of \code{msh\_tjunctions\_tn.ipynb} (section 6), including the overlaps $z_{ik}(t) = \bra{\gamma_i(0)}\psi_k(t)\rangle$ at every snapshot. Not committed (\code{*.jls}); about 2.5 h of TDVP with six threads to create.\\
\code{logical\_gates\_tn.csv} & $M$ and the leakage per gate and stage, from the TN notebook; committed; read by the Python notebook.\\
\code{logical\_gates\_python\_ed\_data/frames.npz} & initial frame $\psi_0$ and the evolved frames after one and two exchanges; not committed; recomputed when missing or with \code{RUN\_EVOLUTION = True}.\\
\bottomrule
\end{tabularx}
\end{center}

\subsection{Results}

\begin{center}
\resizebox{\textwidth}{!}{%
\begin{tabular}{lcccccc}
\toprule
Gate & $P(00\to11)$ & rotation $\theta$ & best ideal gate (overlap) & unexcited prob. & full vs reduced & $\max\bigl||M_{\rm TN}| - |M_{\rm ED}|\bigr|$\\
\midrule
$\sqrt{Z}{}^\dagger$ (one \code{exchange\_12}) & 0 & $90.7180^\circ$ & $\sqrt{Z}{}^\dagger$ (0.999980) & 0.999998 & $3.3\times10^{-12}$ & $4.9\times10^{-9}$\\
$Z$ (two) & 0 & $178.5641^\circ$ & $Z$ (0.999921) & 0.999997 & $1.0\times10^{-11}$ & $8.5\times10^{-9}$\\
$\sqrt{X}{}^\dagger$ (one \code{exchange\_23}) & 0.522078 & $92.5308^\circ$ & $\sqrt{X}{}^\dagger$ (0.999756) & 0.999982 & $9.8\times10^{-8}$ & $1.6\times10^{-8}$\\
$X$ (two) & 0.998051 & $174.9391^\circ$ & $X$ (0.999025) & 0.999964 & $4.2\times10^{-6}$ & $1.8\times10^{-8}$\\
\bottomrule
\end{tabular}}
\end{center}

$\theta$ is \textbf{not a phase}: it is the angle by which the gate rotates the logical Bloch sphere, $U_L\propto e^{-i\theta\,\hat n\cdot\boldsymbol\sigma/2}$
(axes $\hat n = -\hat z$, $+\hat z$, $-\hat x$, $+\hat x$ for the four rows); for the $X$-type gates it sets the amplitudes,
$P(00\to11) = \sin^2\frac\theta2$. The excess over $90^\circ$ (or $180^\circ$) is the hybridization phase $\int\varepsilon\,dt$ of the exchanged pair.
The relative phases are $\varphi_d = \arg\bra{11}\hat U\ket{11} - \arg\bra{00}\hat U\ket{00} = -90.7180^\circ$ ($\sqrt{Z}{}^\dagger$) and $+178.5641^\circ$ ($Z$), and
$\varphi_o = \arg\bra{11}\hat U\ket{00} - \arg\bra{00}\hat U\ket{11} = 0.0000^\circ$ for both $X$-type gates, as for any rotation about $\pm\hat x$
(section 3.2). ``Full vs reduced'' is the largest entry-wise difference of the even block (full amplitudes divided by the square
root of the unexcited probability, global phase removed); the TN and dense (ED) reduced gates agree in $|M|$, $P$, $\theta$ and the leakage.

%======================================================================
\section{Method}
%======================================================================

\subsection{Quasiparticle operators and the Majorana matrix}

With the Nambu spinor $\Psi$ and $H = \tfrac12\Psi^\dagger H_{\rm BdG}\Psi$, the quasiparticle operator of a BdG vector $w$ is
$b(w) = w^\dagger\Psi$. It is \emph{antilinear} in $w$: $b(c\,w) = c^*b(w)$. The Heisenberg evolution is $\hat U^\dagger\Psi\hat U = S_T\Psi$, with
$S_T = \mathcal T e^{-i\int H_{\rm BdG}dt}$ the single-particle propagator, so
\[
\hat U^\dagger b(w)\hat U = b(S_T^\dagger w),\qquad \hat U b(w)\hat U^\dagger = b(S_Tw).
\]
For self-conjugate Majoranas ($C\gamma = \gamma$, $C$ the particle--hole conjugation) $\hat\gamma_j = b(\gamma_j)$ is Hermitian, and
$\hat U\hat\gamma_j\hat U^\dagger = b(S_T\gamma_j) = \sum_i\hat\gamma_i M_{ij}$ with the real matrix $M_{ij} = \bra{\gamma_i}S_T\ket{\gamma_j}$ (real because
$[S_T, C] = 0$). Without leakage out of the zero-mode subspace $M$ is orthogonal; $\|M^TM - 1\|$ measures the leakage.

\subsection{Spin representation and the Jordan--Wigner convention}

Restricted to the zero-mode Fock space, $\hat U$ acts as a unitary $\tilde U$ with $\tilde U\hat\gamma_j\tilde U^\dagger = \sum_i\hat\gamma_iM_{ij}$.
With $M = e^A$ ($A$ real antisymmetric), $\tilde U = \exp(\tfrac14\sum_{ij}A_{ij}\hat\gamma_i\hat\gamma_j)$, up to a global phase and the sign
of $A$ (SO(4) is double covered); the sign is fixed by checking the defining relation.

The $\hat\gamma_i$ are represented on $\ket{n_An_B}$ by Jordan--Wigner, with $f_A = a\otimes1$, $f_B = Z\otimes a$. \textbf{The convention matters}:
$f_m$ must annihilate the positive-energy mode $\psi_m$, so that $n_m = 1$ means $\psi_m$ occupied as in the full method. By antilinearity,
\[
b(\psi_m) = b\Bigl(\tfrac{\gamma_{2m-1} + i\gamma_{2m}}{\sqrt2}\Bigr) = \frac{\hat\gamma_{2m-1} - i\hat\gamma_{2m}}{\sqrt2}
\quad\Rightarrow\quad \hat\gamma_{2m-1} = f_m + f_m^\dagger,\qquad \hat\gamma_{2m} = +i(f_m - f_m^\dagger).
\]
The opposite sign, $\hat\gamma_{2m} = -i(f_m - f_m^\dagger)$, used in the earlier \code{msh\_logical\_gate.jl} / \code{msh\_pfaffian} work,
makes $f_m$ the creation operator of $\psi_m$: it swaps the labels $\ket{00}\leftrightarrow\ket{11}$ and conjugates every relative phase
(one \code{exchange\_12} then appears as $\mathrm{diag}(1, e^{+i\,90.72^\circ})$ instead of $\mathrm{diag}(1, e^{-i\,90.72^\circ})$). The Python
notebook checks the convention against the full Pfaffian amplitudes entry by entry (section~\ref{sec:fullvsred}).

\subsection{Full method: generalized Wick theorem}

A Gaussian state is a set of $2N$ occupied columns of the frame $\psi_0$ (the fixed BdG unitary of \code{fix\_unitary}: $\psi_A$, $\psi_B$ in columns
$2N$, $2N+1$, their partners $C\psi$ in $2N-1$, $2N-2$). For occupations $k$, $\ket k = a^\dagger_{m_1}a^\dagger_{m_2}\cdots\ket{00}$ ($m_1<m_2$),
where the creation operator of mode $m$ is $b(C\phi_m)$ and its adjoint $b(\phi_m)$, $\phi_m$ = column $2N+m$. Writing both states relative to a
reference $\ket r$,
\[
\frac{\bra j\hat U\ket k}{\bra r\hat U\ket r} = \bra r X(r)\,X^\dagger(j)\,\hat U\,X(k)\,X^\dagger(r)\ket r_{\rm normalized},
\]
and moving the operators left of $\hat U$ through it with $S_T^\dagger$, the normalized expectation value of the operator string is the
Pfaffian of the antisymmetric matrix of pairwise contractions
\[
C_{ij} = x_i^\dagger(1-\rho)\,Cx_j,\qquad 1-\rho = 1 - W_rG^{-1}W_L^\dagger,\quad W_L = S_T^\dagger W_r,\quad G = W_L^\dagger W_r,
\]
with $W_r$ the occupied columns of $\ket r$. $1-\rho$ is the mixed (transition) density between $\bra r$ and $\hat U\ket r$; this is where the
whole occupied frame and its evolution enter. $|\bra r\hat U\ket r| = |\det G|^{1/2}$; its sign is not fixed but common to all amplitudes,
so moduli and relative phases are meaningful and the global phase is not. The probability of staying unexcited for input $k$ is
$\sum_j|\bra j\hat U\ket k|^2$ over the logical outputs $j$ (odd outputs vanish by parity).

%======================================================================
\section{Part I --- \code{logical\_gates\_tn.ipynb} (Julia)}
%======================================================================

No time evolution: everything is $4\times4$ algebra on the overlaps cached by \code{msh\_tjunctions\_tn.ipynb}; the notebook runs in seconds.

\subsection{Section 1: setup and the cached overlaps}

\begin{function}{T1}
Duration of one exchange, as \code{duration} in \code{msh\_tjunctions\_tn.ipynb}: $T_1 = (0.1\,(4L_1 + 2L_2 + 8) + 1)\cdot200 = 2040$ for
\code{exchange\_12} and $(0.1\,(L_1 + 3L_2 + 2L_1 + 3) + 1)\cdot400 = 3880$ for \code{exchange\_23} with $L_1 = L_2 = 14$.
\end{function}

\begin{function}{runs, overlaps(kind, t)}
\code{runs[kind]} is the deserialized cache \code{msh\_tjunctions\_tdvp\_\$kind.jls} (\code{kind} = \code{"Z"} or \code{"X"}); a missing cache
stops the notebook with an instruction to run \code{msh\_tjunctions\_tn.ipynb} first. \code{overlaps(kind, t)} returns the $4\times2$ matrix
$z_{ik} = \bra{\gamma_i(0)}\psi_k(t)\rangle$ at the snapshot $t$ (an error if there is none; snapshots are every 5 time units during the gate,
so $T_1$ and $2T_1$ are included). Printed: $T_1 = 2040$ (Z, snapshots to $t = 4460$) and $3880$ (X, to $8140$).
\end{function}

\subsection{Section 2: from the Majorana matrix to the logical gate}

\begin{function}{braid\_matrix(z)}
$M$ from the overlaps. With $\psi_k = (\gamma_{2k-1} + i\gamma_{2k})/\sqrt2$, $C\gamma = \gamma$ and $[S_T, C] = 0$:
$\gamma_{2k-1}(T) = (\psi_k(T) + C\psi_k(T))/\sqrt2$, $\gamma_{2k}(T) = (\psi_k(T) - C\psi_k(T))/(i\sqrt2)$, and
$\bra{\gamma_i}C\psi\rangle = \bra\psi\gamma_i\rangle = z_{ik}^*$, so
\[
M_{i,2k-1} = \sqrt2\,\mathrm{Re}\,z_{ik},\qquad M_{i,2k} = \sqrt2\,\mathrm{Im}\,z_{ik}.
\]
Only $\psi_A(T)$ and $\psi_B(T)$ are needed: the TDVP never evolves the partners, which follow from $C$ exactly.
\end{function}

\begin{function}{jw\_majoranas()}
The four $\hat\gamma_i$ as $4\times4$ matrices on $\ket{n_An_B}$, order $00, 01, 10, 11$: $f_A = a\otimes1$, $f_B = Z\otimes a$ with $a = \ket0\bra1$,
$\hat\gamma_{2m-1} = f_m + f_m^\dagger$, $\hat\gamma_{2m} = +i(f_m - f_m^\dagger)$ (the convention of section 2.2).
\end{function}

\begin{function}{logical\_unitary(M)}
Replaces $M$ by its nearest orthogonal matrix $M_o = UV^T$ (SVD; this renormalizes away the leakage), requires $\det M_o > 0$ (a
parity-preserving rotation), takes $A = \log M_o$ (antisymmetrized), and for both signs of $A$ forms $\tilde U = \exp(\pm\tfrac14\sum A_{ij}\hat\gamma_i\hat\gamma_j)$,
keeping the one with the smaller residual $\max_j\|\tilde U\hat\gamma_j\tilde U^\dagger - \sum_i\hat\gamma_iM_{o,ij}\|$. Returns $(\tilde U, \text{residual})$;
the residuals are $(1$--$3)\times10^{-15}$.
\end{function}

\begin{function}{IDEAL, logical\_even\_block(Ũ), deg(x)}
\code{IDEAL} lists the ideal gates $\sqrt X = e^{-i\pi X/4}$, $\sqrt{X}{}^\dagger$, $X$, $\sqrt Z = \mathrm{diag}(1, i)$, $\sqrt{Z}{}^\dagger$, $Z$.
\code{logical\_even\_block} takes the even block $U_L$ of $\tilde U$ on $\{\ket{00},\ket{11}\}$ (the logical qubit) and describes it in two
complementary ways.
\begin{itemize}
\item \textbf{Rotation angle $\theta$ and axis $\hat n$ (not phases).} With the global phase removed,
$W = U_L/\sqrt{\det U_L} = \cos\frac\theta2 - i\sin\frac\theta2\,\hat n\cdot\boldsymbol\sigma = e^{-i\theta\,\hat n\cdot\boldsymbol\sigma/2}$: the gate rotates the
logical Bloch sphere by $\theta$ about $\hat n$. $\theta = 2\arccos(\mathrm{Re}\,\mathrm{Tr}\,W/2)$, $\hat n_k = \mathrm{Re}(i\,\mathrm{Tr}\,\sigma_kW)/(2\sin\frac\theta2)$.
Ideal values: $\sqrt X$: $\theta = 90^\circ$, $\hat n = \hat x$ ($\sqrt{X}{}^\dagger$: $-\hat x$); $X$: $180^\circ$; $\mathrm{diag}(1, e^{i\varphi})$: $\theta = |\varphi|$,
$\hat n = \mp\hat z$. The amplitudes follow from it, $P(00\to11) = |U_{L,21}|^2 = \sin^2\frac\theta2\,(1 - \hat n_z^2)$ (printed as a check).
$(\theta, \hat n)$ and $(360^\circ - \theta, -\hat n)$ are the same gate up to a global sign; the code reports the pair given by the principal
$\sqrt{\det U_L}$. Two over-rotated $\sqrt{X}{}^\dagger$ exchanges ($2\times92.53^\circ = 185.06^\circ$ about $-\hat x$) therefore appear as
$\theta = 174.94^\circ$ about $+\hat x$.
\item \textbf{Relative phases} $\varphi_d = \arg U_{L,22} - \arg U_{L,11} = \arg\bra{11}\hat U\ket{11} - \arg\bra{00}\hat U\ket{00}$ (the gate phase of a
diagonal gate) and $\varphi_o = \arg U_{L,21} - \arg U_{L,12} = \arg\bra{11}\hat U\ket{00} - \arg\bra{00}\hat U\ket{11}$ (between the off-diagonal elements),
in degrees, \code{NaN} (printed ``--'' by \code{deg}) when the elements vanish. For a rotation about an axis in the $xy$ plane at azimuth
$\phi$, $U_{L,21}/U_{L,12} = e^{2i\phi}$, so $\varphi_o = 2\phi$: $0$ for every rotation about $\pm\hat x$, whatever $\theta$, and $180^\circ$ about $\pm\hat y$.
Both phases are independent of the global phase and of the sign convention of $\ket{11}$.
\end{itemize}
The phase of an off-diagonal element relative to a diagonal one, $\arg\bra{11}\hat U\ket{00} - \arg\bra{00}\hat U\ket{00}$, is $\pm90^\circ$ for a rotation
about $\hat x$, with the sign of $\cos\frac\theta2$: it flips when the rotation passes $180^\circ$ and is undefined for an exact $X$, whose diagonal
vanishes, so it is not used. Finally the function returns the global-phase-insensitive overlaps $|\mathrm{Tr}(V^\dagger U_L)|/2$ with every ideal gate $V$.
\end{function}

\paragraph{Check.} An ideal braid $\gamma_2\to\gamma_3$, $\gamma_3\to-\gamma_2$ gives residual $8\times10^{-16}$, $\theta = 90^\circ$, $\hat n = (-1, 0, 0)$:
exactly $\sqrt{X}{}^\dagger$ (overlap 1.000000).

\subsection{Section 3: logical gates after one and two exchanges}

For $(Z, T_1)$, $(Z, 2T_1)$, $(X, T_1)$, $(X, 2T_1)$: $M$ from \code{braid\_matrix(overlaps(...))}, the leakage $1-\sum_i|z_{ik}|^2$ of each mode,
$\|M^TM - 1\|$, the spin-representation residual, $P(00\to11)$ with the check $\sin^2\frac\theta2(1-\hat n_z^2)$ (equal to six digits), the
rotation $\theta$ about $\hat n$, the relative phases $\varphi_d$, $\varphi_o$ and the ideal-gate overlaps. Printed values:

\begin{center}
\resizebox{\textwidth}{!}{%
\begin{tabular}{lcccccccc}
\toprule
Gate & $t$ & leak $(\psi_A, \psi_B)$ & $\|M^TM-1\|$ & $P(00\to11)$ & rotation $\theta$ & $\hat n$ & $\varphi_d$ & $\varphi_o$\\
\midrule
$\sqrt{Z}{}^\dagger$ & 2040 & $(9.72\times10^{-7}, 1.65\times10^{-11})$ & $1.4\times10^{-6}$ & 0 & $90.7180^\circ$ & $(0, 0, -1)$ & $-90.7180^\circ$ & --\\
$Z$ & 4080 & $(1.99\times10^{-6}, 2.07\times10^{-11})$ & $2.8\times10^{-6}$ & 0 & $178.5641^\circ$ & $(0, 0, 1)$ & $+178.5641^\circ$ & --\\
$\sqrt{X}{}^\dagger$ & 3880 & $(3.49\times10^{-5}, 1.23\times10^{-7})$ & $7.0\times10^{-5}$ & 0.522078 & $92.5308^\circ$ & $(-1, 0, 0)$ & $-0.0002^\circ$ & $0.0000^\circ$\\
$X$ & 7760 & $(3.42\times10^{-5}, 3.49\times10^{-5})$ & $1.0\times10^{-4}$ & 0.998051 & $174.9391^\circ$ & $(1, 0, 0)$ & $+0.0050^\circ$ & $0.0000^\circ$\\
\bottomrule
\end{tabular}}
\end{center}

The ideal-gate overlaps are $0.999980$ ($\sqrt{Z}{}^\dagger$), $0.999921$ ($Z$), $0.999756$ ($\sqrt{X}{}^\dagger$) and $0.999025$ ($X$). One
\code{exchange\_12} leaves pair B untouched (its leakage stays at $10^{-11}$); one \code{exchange\_23} acts on both pairs.

\subsection{Section 4: export}

\code{logical\_gates\_tn.csv}: columns \code{gate, stage, t, leak\_A, leak\_B, M11 \dots M44} (row-major), one row per gate and stage.

%======================================================================
\section{Part II --- \code{logical\_gates\_python\_ed.ipynb} (Python)}
%======================================================================

\subsection{Code from \code{tjunctions\_model.py}}

The notebook imports the model (\code{from tjunctions\_model import *}); these functions are documented in Part II of the T-junction
reference and only summarized here.

\begin{center}
\begin{tabularx}{\textwidth}{lX}
\toprule
Name & Role here\\
\midrule
\code{H0}, \code{ham(f)}, \code{ham\_sp(f)} & $H_{\rm BdG}$ of the initial texture $f_0$; dense and sparse $H$ for a texture $f(c,s,n)$\\
\code{Exchange12}, \code{Exchange23}, \code{f\_array\_12}, \code{f\_array\_23} & the two protocols (durations, textures)\\
\code{zero\_mode\_frame(H)} & \code{eigh}, \code{individual\_MZMs} (Takagi + localization), \code{fix\_unitary}: the fixed frame $\psi_0$ with $\psi_A$, $\psi_B$ in columns $2N$, $2N+1$ and partners in $2N-1$, $2N-2$\\
\code{step(psi, t, dt, ft)}, \code{DT} & 2-stage Gauss--Legendre step (sparse LU), $dt = 0.25$\\
\code{Cop(v)} & particle--hole conjugation of the kwant spinor, $C = \tau_y\sigma_yK$\\
\code{N}, \code{xs} & 144 sites; $x$ coordinate of each site\\
\bottomrule
\end{tabularx}
\end{center}

\subsection{Section 1: setup}

Cell 1 sets \code{OMP\_NUM\_THREADS}, \code{OPENBLAS\_NUM\_THREADS}, \code{MKL\_NUM\_THREADS} to 1 before NumPy is imported (single thread,
as in \code{tjunctions\_python\_ed.ipynb}), imports the model and builds \code{psi0, E0, mbs, mask = zero\_mode\_frame(H0)}; \code{states} lists
the occupations $(0,0), (0,1), (1,0), (1,1)$. Printed: $N = 144$, $D = 576$, $\|\psi_0^\dagger\psi_0 - 1\| = 6.5\times10^{-14}$, mode energies
$1.09\times10^{-9}$ and $1.16\times10^{-9}$.

\subsection{Section 2: time evolution of the full frame}

\begin{function}{T1, gate\_f}
$T_1$ from \code{p12.duration()} and \code{p23.duration()} (\code{Exchange23(rot\_T=400.0)}); \code{gate\_f[kind](t)} is the texture of two
consecutive exchanges, the second one starting at $T_1$.
\end{function}

\begin{function}{evolve\_frame(ft, psi, t0, t1, dt=DT)}
Evolves the whole $4N\times4N$ frame from $t_0$ to $t_1$ with \code{step}, renormalizing the columns after each step: \code{evolve} of
\code{tjunctions\_python\_ed.ipynb} without the readout.
\end{function}

\paragraph{Cache.} If \code{frames.npz} is missing or \code{RUN\_EVOLUTION = True}, both gates are evolved and saved after one and two exchanges
(\code{Z1, Z2, X1, X2}) together with $\psi_0$. When loading, the notebook uses the cached $\psi_0$ (the phases of degenerate bulk eigenvectors
from \code{eigh} depend on the BLAS build and threads, so a recomputed $\psi_0$ need not match bit for bit) and checks that it is an eigenframe of
$H_0$ whose zero-mode columns equal those of \code{zero\_mode\_frame(H0)}. The propagators are $S_T = \psi(T)\psi_0^\dagger$; printed
$\|S_T^\dagger S_T - 1\| = 4\times10^{-12}$ to $1.3\times10^{-11}$. Evolution times of the cached run: Z 692~s, X 1363~s (single thread, measured
while other jobs ran on the machine).

\subsection{Section 3: full method}

\begin{function}{pfaffian(A)}
Pfaffian of a complex antisymmetric matrix by Parlett--Reid elimination with pivoting: columns in pairs $(k, k+1)$; the largest entry of column
$k$ below the diagonal is swapped into row and column $k+1$ (each swap flips the sign); the pivot $A_{k,k+1}$ is a factor; the trailing block gets
the antisymmetric rank-2 update $A \leftarrow A + \tau a^T - a\tau^T$, $\tau = A_{k,k+2:}/A_{k,k+1}$. Check: $|\mathrm{Pf}(A)^2 - \det A| = 3.7\times10^{-11}$
for a random $8\times8$ matrix.
\end{function}

\begin{function}{Cvec, mask\_of(bits)}
\code{Cvec = Cop}. \code{mask\_of} marks columns $2N + m$ with $n_m = 1$ and mirrors: column $2N-1-m$ is occupied exactly when $2N+m$ is empty, so
every state occupies $2N$ columns (the CZ notebooks' \code{masks[:, :2N] = flip(\textasciitilde masks[:, 2N:])}).
\end{function}

\begin{function}{signed\_amplitudes(pairs, S\_T, psi0, candidates=states)}
$\bra j\hat U\ket k$ for all $(j, k)$ in \code{pairs}, up to one common sign (section 2.3). The reference $r$ is the candidate with the largest
$|\det(W_r^\dagger S_TW_r)|$; $W_L = S_T^\dagger W_r$, $G = W_L^\dagger W_r$, \code{norm\_r} $= |\det G|^{1/2}$. For each pair the operator vectors are
$S_T^\dagger X(r)$, $S_T^\dagger X^\dagger(j)$, $X(k)$, $X^\dagger(r)$ with $X(\text{bits}) = [C\phi_m]$ for occupied $m$ (ascending) and $X^\dagger = [\phi_m]$
(descending); the contraction matrix is $P^\dagger CP - (P^\dagger W_r)G^{-1}(W_L^\dagger CP)$ (the $4N\times4N$ mixed density is never formed),
its strict upper triangle antisymmetrized, and the amplitude is its Pfaffian times \code{norm\_r}. An empty string gives \code{norm\_r}.
Returns the amplitudes, $r$ and \code{norm\_r}. This is \code{projected\_amplitudes} of \code{period\_finding\_signed\_amplitudes.py} without
mid-protocol projections.
\end{function}

\paragraph{Checks and references.} With $S_T = 1$, $\max|A_{jk} - \delta_{jk}| = 3.9\times10^{-15}$. References chosen:
$r = (0,1)$ with $|\bra r\hat U\ket r| = 0.999999$ for Z1 and Z2, $0.691313$ for X1, and $r = (0,0)$ with $0.044149$ for X2 (an $X$ gate empties the
diagonal; the amplitudes are still consistent, see below).

\subsection{Section 4: reduced method}

\begin{function}{gamma0, braid\_matrix(S\_T)}
\code{gamma0} holds $\gamma_{2k-1} = (\psi_k + C\psi_k)/\sqrt2$ and $\gamma_{2k} = (\psi_k - C\psi_k)/(i\sqrt2)$ from the frame (the same convention as the
TN notebook); printed $\|\gamma^\dagger\gamma - 1\| = 2.3\times10^{-15}$, $\max\|C\gamma - \gamma\| = 0$, centres $x = -46.6, -17.5, 16.5, 45.6$.
\code{braid\_matrix} returns $\mathrm{Re}\,(\gamma^\dagger S_T\gamma)$ and $\max|\mathrm{Im}|$ ($\le 2.3\times10^{-15}$).
\end{function}

\begin{function}{jw\_majoranas(), logical\_unitary(M), IDEAL, phases(UL), logical\_even\_block(U), deg}
The functions of Part I in Python (\code{phases} computes $\varphi_d$, $\varphi_o$ of any $2\times2$ even block, used for both methods) (\code{scipy.linalg.logm}, \code{expm}, \code{svd}), with the same Jordan--Wigner convention. Residuals
$\le 2.4\times10^{-15}$; $\|M^TM - 1\|$ as in the TN notebook.
\end{function}

\subsection{Section 5: full vs reduced}\label{sec:fullvsred}

Per gate and stage: the even-block probabilities of the full method, the unexcited probability per input, the full probabilities
renormalized by it, the reduced probabilities, the relative phases $\varphi_d$ and $\varphi_o$ (section 3.2) of both methods, and
the whole even block entry by entry: $\max|U_{\rm full}/\sqrt{P_{\rm stay}} - \tilde U|$ after removing the global phase (fixed on the largest
entry).

\begin{center}
\resizebox{\textwidth}{!}{%
\begin{tabular}{lcccc}
\toprule
Gate & unexcited prob. & $\varphi_d = \arg\bra{11}\hat U\ket{11} - \arg\bra{00}\hat U\ket{00}$ & $\varphi_o = \arg\bra{11}\hat U\ket{00} - \arg\bra{00}\hat U\ket{11}$ & even block\\
 & & full / reduced & full / reduced & entry by entry\\
\midrule
$\sqrt{Z}{}^\dagger$ & 0.999998 & $-90.7180^\circ$ / $-90.7180^\circ$ & -- & $3.3\times10^{-12}$\\
$Z$ & 0.999997 & $+178.5641^\circ$ / $+178.5641^\circ$ & -- & $1.0\times10^{-11}$\\
$\sqrt{X}{}^\dagger$ & 0.999982 & $-0.0002^\circ$ / $-0.0002^\circ$ & $0.0000^\circ$ / $0.0000^\circ$ & $9.8\times10^{-8}$\\
$X$ & 0.999964 & $+0.0050^\circ$ / $+0.0050^\circ$ & $0.0000^\circ$ / $0.0000^\circ$ & $4.2\times10^{-6}$\\
\bottomrule
\end{tabular}}
\end{center}

The renormalized full probabilities equal the reduced ones to $3.5\times10^{-29}$ (Z) and $1$--$3\times10^{-11}$ (X). The remaining entry-wise
differences of the $X$-type gates ($10^{-7}$--$10^{-6}$) are of the order of the bulk excitation ($1 - P_{\rm stay} = 1.8\times10^{-5}$ and
$3.6\times10^{-5}$), which only the full method contains. With the opposite Jordan--Wigner sign the phases come out conjugated and the entry-wise
difference is of order 1. Both methods give $\varphi_o = 0$ to the printed four decimals: the $X$-type gates are rotations about $\pm\hat x$, the
departure from the ideal gates is only in the rotation angle ($92.53^\circ$ instead of $90^\circ$).

\subsection{Section 6: comparison with the tensor-network gates}

Reads \code{logical\_gates\_tn.csv}, rebuilds $\tilde U$ from $M_{\rm TN}$ with the Python \code{logical\_unitary}, and compares quantities that do
not depend on the sign of each $\gamma$ (fixed independently in the two codes): $\max\bigl||M_{\rm TN}| - |M_{\rm ED}|\bigr|$, $P(00\to11)$, $\theta$ and the
leakage $1 - \sum_i|\bra{\gamma_i}S_T\psi_k\rangle|^2$.

\begin{center}
\resizebox{\textwidth}{!}{%
\begin{tabular}{lccccc}
\toprule
Gate & $\max\bigl||M_{\rm TN}|-|M_{\rm ED}|\bigr|$ & $P(00\to11)$ TN / ED & $\theta$ TN / ED & leak TN $(A, B)$ & leak ED $(A, B)$\\
\midrule
$\sqrt{Z}{}^\dagger$ & $4.9\times10^{-9}$ & 0 / 0 & $90.7180^\circ$ / $90.7180^\circ$ & $(9.72\times10^{-7}, 1.7\times10^{-11})$ & $(9.74\times10^{-7}, 0)$\\
$Z$ & $8.5\times10^{-9}$ & 0 / 0 & $178.5641^\circ$ / $178.5641^\circ$ & $(1.99\times10^{-6}, 2.1\times10^{-11})$ & $(2.00\times10^{-6}, 0)$\\
$\sqrt{X}{}^\dagger$ & $1.6\times10^{-8}$ & 0.522078 / 0.522078 & $92.5308^\circ$ / $92.5308^\circ$ & $(3.49\times10^{-5}, 1.23\times10^{-7})$ & $(3.49\times10^{-5}, 1.24\times10^{-7})$\\
$X$ & $1.8\times10^{-8}$ & 0.998051 / 0.998051 & $174.9391^\circ$ / $174.9391^\circ$ & $(3.42\times10^{-5}, 3.49\times10^{-5})$ & $(3.42\times10^{-5}, 3.49\times10^{-5})$\\
\bottomrule
\end{tabular}}
\end{center}

The MPS/TDVP evolution of the two mode vectors (bond dimension $\le 64$, $dt = 0.1$) and the dense Gauss--Legendre evolution of the full
frame ($dt = 0.25$) give the same logical gates to $\sim10^{-8}$; the TN route needs only the two vectors and no dense matrix.

%======================================================================
\section{\code{tjunctions\_model.py}}
%======================================================================

The Python model of the array, moved out of \code{tjunctions\_python\_ed.ipynb} (whose outputs are unchanged by the move) so that both Python
notebooks use one copy. Contents, in order: environment variables for single-thread BLAS (effective only if the module is imported before
NumPy; the notebooks also set them in their first cell); Pauli matrices and their Kronecker products; geometry (\code{NLEG}, \code{NCELL},
\code{center\_x}, \code{array\_pos}, \code{label}, \code{build\_system\_3T}, and the module-level \code{sys}, \code{N}, \code{tags}, \code{xs}, \code{ys});
textures and protocols (\code{smoothstep}, \code{f0}, \code{spins}, \code{Exchange12}, \code{Exchange23}, \code{Fusion} and their texture functions);
the Hamiltonian (\code{params}, \code{ham\_kwant}, \code{H\_noJ}, \code{exchange\_matrix}, \code{ham}, \code{Cop}, \code{H0}); the zero-mode
construction (\code{dag}, \code{PHconj}, \code{takagi}, \code{individual\_MZMs}, \code{fix\_unitary}); two functions added with the move,
\code{zero\_mode\_frame(H)} (returns \code{psi0, E, mbs, mask}) and \code{logical\_masks()} (the $4\times4N$ masks of $\ket{00}$, $\ket{01}$, $\ket{10}$,
$\ket{11}$); the sparse Hamiltonian and Gauss--Legendre step (\code{exchange\_sparse}, \code{ham\_sp}, \code{step}, \code{DT}); and
\code{electron\_density(W)}. All are described in Part II of \code{docs/msh\_tjunctions\_function\_reference.pdf}.

\end{document}
